% \begin{document}
\chapter{Polymorphism}

\section{Parametric Polymorphism}

In the simply-typed $\lambda$-calculus every argument carries a type annotation,
so each type needs its own identity function: $\lambda x : \kw{int}.\, x$,
$\lambda x : \kw{bool}.\, x$, $\lambda x : \kw{int} \to \kw{int}.\, x$, and so on.
\emph{Polymorphism} is the ability to use code with values of different types.
Three common kinds:
\begin{itemize}
  \item \emph{Subtype polymorphism}: a subsumption rule lets a value of type
        $\tau$ be used as type $\tau'$ whenever $\tau$ is a subtype of $\tau'$.
  \item \emph{Ad-hoc polymorphism}: code that looks polymorphic but is not
        (e.g.\ overloading); the argument types dispatch, at compile time or run
        time, to one of several separate implementations.
  \item \emph{Parametric polymorphism}: code written without knowledge of the
        actual argument types (polymorphic functions in OCaml; generics in Java,
        Scala, Rust). This is the subject of this lecture.
\end{itemize}

\section{System F}

System F extends the simply-typed $\lambda$-calculus with abstraction over types.
It was invented independently by Girard (1972, ``System F'') and Reynolds (1974,
``polymorphic $\lambda$-calculus''). A \emph{type abstraction} $\Lambda X.\, e$ is
a function from types to terms; a \emph{type application} $e\,[\tau]$ instantiates
it at $\tau$. For example, the polymorphic identity is
\[
\mathit{ID} \triangleq \Lambda X.\ \lambda x : X.\ x.
\]

\subsection{Syntax and operational semantics}
\begin{align*}
e &::= x \mid \lambda x : \tau.\, e \mid e_1\, e_2 \mid () \mid n \mid e_1 \oplus e_2
      \mid \ttkw{true} \mid \ttkw{false} \mid \ttkw{if}\ e_1\ \ttkw{then}\ e_2\ \ttkw{else}\ e_3
      \mid \Lambda X.\, e \mid e\,[\tau] \\
v &::= \lambda x : \tau.\, e \mid () \mid n \mid \ttkw{true} \mid \ttkw{false} \mid \Lambda X.\, e \\
\tau &::= \kw{unit} \mid \kw{int} \mid \kw{bool} \mid \tau_1 \to \tau_2 \mid X \mid \forall X.\, \tau
\end{align*}
Here $\oplus \in \{+, -, \times\}$. Both $\Lambda X.\, e$ and $\forall X.\, \tau$ bind
$X$, giving the usual free/bound type variables and $\alpha$-equivalence (e.g.\
$\forall X.\, X \to X$ and $\forall Y.\, Y \to Y$ are the same type). Evaluation is
CBV as before, with one new context and one new rule:
\[
E ::= \cdots \mid E\,[\tau]
\qquad\qquad
\inferrule*[Right=\textsc{Type-$\beta$}]
  { }
  {(\Lambda X.\, e)\,[\tau] \to \tsub{e}{\tau}{X}}
\]
Note $\Lambda X.\, e$ is a value: we do not evaluate under $\Lambda$ (just as not
under $\lambda$). For example,
\[
\mathit{ID}\,[\kw{int}]\ 5 = (\Lambda X.\, \lambda x : X.\, x)\,[\kw{int}]\ 5
  \to (\lambda x : \kw{int}.\, x)\ 5 \to 5,
\]
and $\mathit{ID}\,[\kw{int} \to \kw{int}]$ steps to $\lambda x : \kw{int} \to \kw{int}.\, x$.

Type substitution must be capture-avoiding. E.g.\ in
$\Lambda Y.\, (\Lambda X.\, \Lambda Y.\, \lambda x : X.\, \lambda y : Y.\, x)\,[Y]$
we must rename the inner bound variable:
\[
(\forall Y.\, X \to Y \to X)\{Y/X\} = \forall Y'.\, Y \to Y' \to Y,
\]
not $\forall Y.\, Y \to Y \to Y$, which would capture the outer $Y$.

\subsection{Type system}

$\Lambda X.\, e$ has type $\forall X.\, \tau$ where $\tau$ is the type of $e$
(which may mention $X$). Typing now tracks, besides the term context $\Gamma$, a
set $\Delta$ of type variables in scope, so judgments are
$\Delta, \Gamma \vdash e : \tau$; e.g.\ $\lambda x : X.\, x$ is ill-typed on its
own. An auxiliary judgment $\Delta \vdash \tau\ \mathit{ok}$ says every free type
variable of $\tau$ is in $\Delta$. The rules of the simply-typed fragment carry
$\Delta$ unchanged, and annotations must be well-formed:
\[
\inferrule*[Right=\textsc{T-Var}]
  {\Gamma(x) = \tau}
  {\Delta, \Gamma \vdash x : \tau}
\qquad\qquad
\inferrule*[Right=\textsc{T-Abs}]
  {\Delta \vdash \tau\ \mathit{ok} \\ \Delta, \Gamma, x : \tau \vdash e : \tau'}
  {\Delta, \Gamma \vdash \lambda x : \tau.\, e : \tau \to \tau'}
\]
\[
\inferrule*[Right=\textsc{T-App}]
  {\Delta, \Gamma \vdash e_1 : \tau \to \tau' \\ \Delta, \Gamma \vdash e_2 : \tau}
  {\Delta, \Gamma \vdash e_1\, e_2 : \tau'}
\qquad\qquad
\inferrule*[Right=\textsc{T-If}]
  {\Delta, \Gamma \vdash e_1 : \kw{bool} \\ \Delta, \Gamma \vdash e_2 : \tau \\ \Delta, \Gamma \vdash e_3 : \tau}
  {\Delta, \Gamma \vdash \ttkw{if}\ e_1\ \ttkw{then}\ e_2\ \ttkw{else}\ e_3 : \tau}
\]
Rules for $()$, integers, arithmetic, and Booleans are similar. The two new rules:
\[
\inferrule*[Right=\textsc{T-TAbs}]
  {\Delta \cup \{X\}, \Gamma \vdash e : \tau}
  {\Delta, \Gamma \vdash \Lambda X.\, e : \forall X.\, \tau}
\qquad\qquad\quad
\inferrule*[Right=\textsc{T-TApp}]
  {\Delta, \Gamma \vdash e : \forall X.\, \tau' \\ \Delta \vdash \tau\ \mathit{ok}}
  {\Delta, \Gamma \vdash e\,[\tau] : \tsub{\tau'}{\tau}{X}}
\]
In \textsc{T-TAbs} we choose $X \notin \Delta$ by $\alpha$-renaming. Well-formedness
of types:
\[
\inferrule{X \in \Delta}{\Delta \vdash X\ \mathit{ok}}
\quad
\Delta \vdash \kw{unit}\ \mathit{ok}
\quad
\Delta \vdash \kw{int}\ \mathit{ok}
\quad
\Delta \vdash \kw{bool}\ \mathit{ok}
\quad
\inferrule{\Delta \vdash \tau_1\ \mathit{ok} \\ \Delta \vdash \tau_2\ \mathit{ok}}{\Delta \vdash \tau_1 \to \tau_2\ \mathit{ok}}
\quad
\inferrule{\Delta \cup \{X\} \vdash \tau\ \mathit{ok}}{\Delta \vdash \forall X.\, \tau\ \mathit{ok}}
\]
Derivation that the polymorphic identity has type $\forall X.\, X \to X$:
\[
\inferrule*[Right=\textsc{T-TAbs}]
{
  \inferrule*[Right=\textsc{T-Abs}]
  {
    \inferrule{X \in \{X\}}{\{X\} \vdash X\ \mathit{ok}}
    \\
    \inferrule*[Right=\textsc{T-Var}]
      {(x : X)(x) = X}
      {\{X\}, x : X \vdash x : X}
  }
  {\{X\}, \cdot \vdash \lambda x : X.\, x : X \to X}
}
{\emptyset, \cdot \vdash \Lambda X.\, \lambda x : X.\, x : \forall X.\, X \to X}
\]
By \textsc{T-TApp}, $\mathit{ID}\,[\kw{int}]$ has type
$(X \to X)\{\kw{int}/X\} = \kw{int} \to \kw{int}$.

\section{Programming in System F}

\subsection{Self application}
The term $\lambda x.\, x\,x$ is untypable in the simply-typed $\lambda$-calculus,
but in System F we can give $x$ a polymorphic type and instantiate it at that same
type:
\[
\vdash \lambda x : \forall X.\, X \to X.\ x\,[\forall X.\, X \to X]\ x
  \ :\ (\forall X.\, X \to X) \to (\forall X.\, X \to X).
\]
Applied to $\mathit{ID}$ it returns $\mathit{ID}$. We still cannot type $\Omega$,
since every well-typed System F program terminates (Section~\ref{sec:meta}).

\subsection{Church numerals}
All Church numerals share one type:
\[
\mathit{NAT} \triangleq \forall X.\, (X \to X) \to X \to X.
\]
\begin{align*}
0 &\triangleq \Lambda X.\ \lambda f : X \to X.\ \lambda x : X.\ x \\
1 &\triangleq \Lambda X.\ \lambda f : X \to X.\ \lambda x : X.\ f\,x \\
\mathit{SUCC} &\triangleq \lambda n : \mathit{NAT}.\ \Lambda X.\ \lambda f : X \to X.\ \lambda x : X.\ f\,(n\,[X]\ f\ x)
\end{align*}
$\mathit{SUCC}$ must instantiate $n$ at $X$ before applying it to $f$ and $x$. Note
that $\mathit{twice}$ is exactly $2$.

\subsection{Products and sums}
Products and sums are encoded with no new types. A pair waits for a consumer $p$
and hands it both components, with the consumer choosing the result type $R$:
\begin{align*}
\tau_1 \times \tau_2 &\triangleq \forall R.\, (\tau_1 \to \tau_2 \to R) \to R \\
\mathit{pair} &\triangleq \Lambda T_1.\, \Lambda T_2.\ \lambda v_1 : T_1.\ \lambda v_2 : T_2.\
   \Lambda R.\ \lambda p : T_1 \to T_2 \to R.\ p\ v_1\ v_2 \\
\mathit{fst} &\triangleq \Lambda T_1.\, \Lambda T_2.\ \lambda v : T_1 \times T_2.\ v\,[T_1]\ (\lambda x : T_1.\, \lambda y : T_2.\, x) \\
\mathit{snd} &\triangleq \Lambda T_1.\, \Lambda T_2.\ \lambda v : T_1 \times T_2.\ v\,[T_2]\ (\lambda x : T_1.\, \lambda y : T_2.\, y)
\end{align*}
A sum value waits for one handler per case and calls the appropriate one:
\begin{align*}
\tau_1 + \tau_2 &\triangleq \forall R.\, (\tau_1 \to R) \to (\tau_2 \to R) \to R \\
\mathit{inl} &\triangleq \Lambda T_1.\, \Lambda T_2.\ \lambda v : T_1.\ \Lambda R.\ \lambda b_1 : T_1 \to R.\ \lambda b_2 : T_2 \to R.\ b_1\, v \\
\mathit{inr} &\triangleq \Lambda T_1.\, \Lambda T_2.\ \lambda v : T_2.\ \Lambda R.\ \lambda b_1 : T_1 \to R.\ \lambda b_2 : T_2 \to R.\ b_2\, v \\
\mathit{case} &\triangleq \Lambda T_1.\, \Lambda T_2.\, \Lambda R.\ \lambda v : T_1 + T_2.\ \lambda b_1 : T_1 \to R.\ \lambda b_2 : T_2 \to R.\ v\,[R]\ b_1\ b_2
\end{align*}
Even $\kw{unit}$ and $\kw{bool}$ are unnecessary: $\kw{unit} \triangleq \forall R.\, R \to R$
(whose only value is $\mathit{ID}$) and $\kw{bool} \triangleq \kw{unit} + \kw{unit}$.

\section{Existential Types}

A universal type $\forall X.\, \tau$ means the \emph{user} chooses $X$. Dually, an
existential type $\exists X.\, \tau$ means there is \emph{some} type $X$ with a
value of type $\tau$: the \emph{creator} chooses $X$ and the user does not know it.
This hides an implementation behind an interface, as modules do in OCaml. For
example, an interface for counters (an initial counter, an increment, and a read):
\[
\mathit{COUNTER} \triangleq \exists C.\ C \times (C \to C) \times (C \to \kw{int}).
\]
(We write $(e_1, e_2, e_3)$ for tuples and $\#i\ e$ for projections; these may be
added directly or encoded as above.)

\subsection{Syntax and semantics}
\begin{align*}
e &::= \cdots \mid \ttkw{pack}\ \{\tau_1, e\}\ \ttkw{as}\ \exists X.\, \tau_2
      \mid \ttkw{unpack}\ \{X, x\} = e_1\ \ttkw{in}\ e_2 \\
v &::= \cdots \mid \ttkw{pack}\ \{\tau_1, v\}\ \ttkw{as}\ \exists X.\, \tau_2 \\
\tau &::= \cdots \mid \exists X.\, \tau
\end{align*}
$\ttkw{pack}\ \{\tau_1, e\}\ \ttkw{as}\ \exists X.\, \tau_2$ pairs a witness type
$\tau_1$ with a value of type $\tsub{\tau_2}{\tau_1}{X}$.
$\ttkw{unpack}\ \{X, x\} = e_1\ \ttkw{in}\ e_2$ opens the package, binding the hidden
type to $X$ and the value to $x$ in $e_2$ (so $\exists X.\, \tau$ binds $X$ in
$\tau$, and $\ttkw{unpack}$ binds $X$ and $x$ in $e_2$). Semantics:
\[
E ::= \cdots \mid \ttkw{pack}\ \{\tau_1, E\}\ \ttkw{as}\ \exists X.\, \tau_2
  \mid \ttkw{unpack}\ \{X, x\} = E\ \ttkw{in}\ e
\]
\[
\ttkw{unpack}\ \{X, x\} = (\ttkw{pack}\ \{\tau_1, v\}\ \ttkw{as}\ \exists Y.\, \tau_2)\ \ttkw{in}\ e
  \ \to\ \tsub{\tsub{e}{v}{x}}{\tau_1}{X}
\]

\subsection{Typing}
\[
\inferrule*[Right=\textsc{T-Pack}]
  {\Delta \vdash \tau_1\ \mathit{ok} \\ \Delta, \Gamma \vdash e : \tsub{\tau_2}{\tau_1}{X}}
  {\Delta, \Gamma \vdash \ttkw{pack}\ \{\tau_1, e\}\ \ttkw{as}\ \exists X.\, \tau_2 : \exists X.\, \tau_2}
\]
\[
\inferrule*[Right=\textsc{T-Unpack}]
  {\Delta, \Gamma \vdash e_1 : \exists X.\, \tau_1 \\
   \Delta \cup \{X\}, \Gamma, x : \tau_1 \vdash e_2 : \tau_2 \\
   \Delta \vdash \tau_2\ \mathit{ok}}
  {\Delta, \Gamma \vdash \ttkw{unpack}\ \{X, x\} = e_1\ \ttkw{in}\ e_2 : \tau_2}
\]
In \textsc{T-Pack} the annotation $\exists X.\, \tau_2$ says which occurrences of
the witness are hidden (otherwise a value of type $\kw{int} \to \kw{int}$ could be
typed $C \to C$, $C \to \kw{int}$, $\kw{int} \to C$, or $\kw{int} \to \kw{int}$). In
\textsc{T-Unpack}, $e_2$ is checked with $X$ an opaque new type variable ($X \notin
\Delta$), and $\Delta \vdash \tau_2\ \mathit{ok}$ ensures the hidden type does not
escape through the result type.

\subsection{Example: counters}
\[
\mathit{counter} \triangleq \ttkw{pack}\ \{\kw{int},\ (0,\ \lambda i : \kw{int}.\, i + 1,\ \lambda i : \kw{int}.\, i)\}\ \ttkw{as}\ \exists C.\ C \times (C \to C) \times (C \to \kw{int})
\]
A client:
\[
\ttkw{unpack}\ \{C, c\} = \mathit{counter}\ \ttkw{in}\ (\#3\ c)\ ((\#2\ c)\ ((\#2\ c)\ (\#1\ c))) \to^* 2.
\]
Inside the $\ttkw{unpack}$, the client only knows
$c : C \times (C \to C) \times (C \to \kw{int})$; e.g.\ $(\#1\ c) + 1$ is ill-typed,
and $\ttkw{unpack}\ \{C, c\} = \mathit{counter}\ \ttkw{in}\ \#1\ c$ is rejected by
\textsc{T-Unpack} because its type $C$ would escape. Hence implementations are
interchangeable, e.g.\ one that counts down internally:
\[
\mathit{counter}' \triangleq \ttkw{pack}\ \{\kw{int},\ (0,\ \lambda i : \kw{int}.\, i - 1,\ \lambda i : \kw{int}.\, 0 - i)\}\ \ttkw{as}\ \exists C.\ C \times (C \to C) \times (C \to \kw{int}).
\]
No client can distinguish $\mathit{counter}$ from $\mathit{counter}'$ (\emph{representation
independence}, a consequence of parametricity).

\subsection{Encoding existentials}
Existentials are definable from universals: a package takes a result type $Y$ and a
client that works for any hidden type, and calls it with its witness and value.
\begin{align*}
\exists X.\, \tau &\triangleq \forall Y.\, (\forall X.\, \tau \to Y) \to Y \\
\ttkw{pack}\ \{\tau_1, e\}\ \ttkw{as}\ \exists X.\, \tau_2 &\triangleq \Lambda Y.\ \lambda f : \forall X.\, \tau_2 \to Y.\ f\,[\tau_1]\ e \\
\ttkw{unpack}\ \{X, x\} = e_1\ \ttkw{in}\ e_2 &\triangleq e_1\,[\tau_2]\ (\Lambda X.\, \lambda x : \tau_1.\, e_2)
\end{align*}
where $e_1 : \exists X.\, \tau_1$ and $e_2 : \tau_2$. The side condition that $X$ not
occur free in $\tau_2$ appears here as well: the client $\Lambda X.\, \lambda x :
\tau_1.\, e_2$ must have type $\forall X.\, \tau_1 \to \tau_2$, which only makes sense
if $\tau_2$ does not mention $X$.

\section{Metatheory}\label{sec:meta}

System F is type sound, with the same proof structure as for the simply-typed
$\lambda$-calculus (progress and preservation).

\thm{Type soundness}{
If $\emptyset, \cdot \vdash e : \tau$ and $e \to^* e'$ and $e' \not\to$, then $e'$ is
a value.
}

The canonical forms lemma gains a case: if $\emptyset, \cdot \vdash v : \forall X.\,
\tau$ then $v = \Lambda X.\, e$ for some $e$. Preservation for \textsc{Type-$\beta$}
needs a new lemma:

\mlenma{Type substitution}{
If $\Delta \cup \{X\}, \Gamma \vdash e : \tau$ and $\Delta \vdash \tau'\ \mathit{ok}$,
then $\Delta, \tsub{\Gamma}{\tau'}{X} \vdash \tsub{e}{\tau'}{X} : \tsub{\tau}{\tau'}{X}$.
}

\thm{Strong normalization (Girard, 1972)}{
If $\emptyset, \cdot \vdash e : \tau$, then every reduction sequence starting from $e$
is finite.
}

The proof is much harder than for the simply-typed case, since the type of
$\Lambda X.\, e$ quantifies over all types, including $\forall X.\, \tau$ itself;
Girard's proof uses \emph{reducibility candidates}. Consequently System F is not
Turing complete, and $\Omega$, $Y$, and $Z$ remain untypable. It is nonetheless very
expressive: the functions on naturals definable in System F are exactly those
provably total in second-order arithmetic, far beyond the primitive recursive
functions.

\section{Type Erasure}

Types never affect the result of a program, so they can be erased before running:
\begin{align*}
\erase(x) &= x \\
\erase(\lambda x : \tau.\, e) &= \lambda x.\, \erase(e) \\
\erase(e_1\, e_2) &= \erase(e_1)\ \erase(e_2) \\
\erase(\Lambda X.\, e) &= \lambda z.\, \erase(e) \qquad \text{where } z \text{ is fresh for } e \\
\erase(e\,[\tau]) &= \erase(e)\ ()
\end{align*}
The remaining cases erase types in subexpressions. A type abstraction becomes a
function of a dummy argument (so it is still a value), and a type application passes
$()$, so each \textsc{Type-$\beta$} step becomes an ordinary $\beta$ step:
\[
\erase((\Lambda X.\, e)\,[\tau]) = (\lambda z.\, \erase(e))\ () \to \erase(e) = \erase(\tsub{e}{\tau}{X}).
\]

\thm{Erasure}{
For all expressions $e$ and $e'$, $e \to e'$ if and only if $\erase(e) \to \erase(e')$.
}

So we can type check with full types and run without them, as OCaml does.

\section{Parametricity}

A function of type $\forall X.\, X \to X$ knows nothing about $X$: it cannot inspect,
create, or modify a value of type $X$, so it must be the identity. Likewise a
function of type $\forall X.\, X \to X \to X$ can only return one of its arguments,
so there are exactly two (up to equivalence): the Church Booleans. And there is no
closed term of type $\forall X.\, \forall Y.\, X \to Y$, since no value of an
arbitrary type $Y$ can be produced.

This is \emph{parametricity}: Reynolds (1983) showed polymorphic functions treat
type arguments uniformly, and Wadler (1989) showed how to derive ``theorems for
free'' from types. For instance, any $f : \forall X.\, X\ \mathit{list} \to X\
\mathit{list}$ can only rearrange, drop, or duplicate elements, hence for every $g$,
\[
\mathit{map}\ g\ (f\ l) = f\ (\mathit{map}\ g\ l).
\]
% \end{document}
